\section*{第 \S\arabic{exercisegroup} 节习题}
\phantomsection
\addcontentsline{toc}{section}{第 \protect\S\arabic{exercisegroup} 节习题}

\begin{enumerate}
\def\labelenumi{\arabic{enumi})}
\tightlist
\item
  设 \(f\) 是定义在区间 \(I \subset \mathbf{R}\) 上的实变向量函数，在内点 \(x_0\)（\(I\) 的）处可导。证明商
\end{enumerate}

\[
\frac {f (x _ {0} + h) - f (x _ {0} - k)}{h + k}
\]

趋于 \(f'(x_0)\)（当 \(h\) 与 \(k\) 保持取值 \(> 0\) 而趋于 0 时）。逆命题。

*证明函数 \(f\)——等于 \(x^2 \sin 1/x\)（对 \(x \neq 0\)）且等于 0（对 \(x = 0\)）——处处可导，但 \((f(y) - f(z)) / (y - z)\) 并不趋于 \(f'(0)\)（当 \(y\) 与 \(z\) 在保持互异且 \(>0_*\) 的条件下趋于 0 时）。

\begin{enumerate}
\def\labelenumi{\arabic{enumi})}
\setcounter{enumi}{1}
\tightlist
\item
  在区间 I = {[}0, 1{]} 上，我们归纳地定义一个连续实函数序列 \((f_{n})\) 如下：取 \(f_{0}(x) = x\)；对每个整数 \(n \geqslant 1\)，函数 \(f_{n}\) 在 \(3^{n}\) 个区间 \(\left[\frac{k}{3^{n}}, \frac{k+1}{3^{n}}\right]\)（对于 \(0 \leqslant k \leqslant 3^{n}-1\)）的每一个上都是仿射线性的；此外，取
\end{enumerate}

\[
f _ {n + 1} \left(\frac {k}{3 ^ {n}}\right) = f _ {n} \left(\frac {k}{3 ^ {n}}\right)
\]

\[
f _ {n + 1} \left(\frac {k}{3 ^ {n}} + \frac {1}{3 ^ {n + 1}}\right) = f _ {n} \left(\frac {k}{3 ^ {n}} + \frac {2}{3 ^ {n + 1}}\right), f _ {n + 1} \left(\frac {k}{3 ^ {n}} + \frac {2}{3 ^ {n + 1}}\right) = f _ {n} \left(\frac {k}{3 ^ {n}} + \frac {1}{3 ^ {n + 1}}\right).
\]

证明序列 \((f_{n})\) 在 I 上一致收敛于一个连续函数，它在区间 {]}0, 1{[} 的任何点上都没有导数（有限的或无限的）（利用习题 1）。

\begin{enumerate}
\def\labelenumi{\arabic{enumi})}
\setcounter{enumi}{2}
\item
  设 \(\mathcal{C}(\mathrm{I})\) 是定义在紧区间 \(\mathrm{I} = [a,b]\)（含于 \(\mathbf{R}\)）上的连续实值有限函数所成的完备空间，并给 \(\mathcal{C}(\mathrm{I})\) 赋予一致收敛拓扑（Gen.~Top., X, p.~277）。设 A 是 \(\mathcal{C}(\mathrm{I})\) 中由满足如下条件的函数 \(x\) 组成的子集：至少存在一点 \(t\in [a,b[\)（依赖于函数 \(x\)），使函数 \(x\) 在该点有有限的右导数。证明 A 是 \(\mathcal{C}(\mathrm{I})\) 中的瘦集（Gen.~Top., IX, p.~192），因而其余集，即在 \([a,b[\) 的任何一点都没有有限右导数的连续函数所成之集，是 \(\mathcal{C}(\mathrm{I})\) 的一个贝尔子空间（Gen.~Top., IX, p.~192）。（设 \(\mathrm{A}_n\) 是由这样的函数 \(x\in \mathcal{C}(\mathrm{I})\) 组成的集：至少存在 \(t\) 的一个值，满足 \(a\leqslant t\leqslant b - 1 / n\)（且依赖于 \(x\)），使得 \(|x(t') - x(t)|\leqslant n|t' - t|\) 对 \(t'\) 的每个满足 \(t\leqslant t'\leqslant t + 1 / n\) 的值都成立。证明每个 \(\mathrm{A}_n\) 都是 \(\mathcal{C}(\mathrm{I})\) 中的闭的无处稠密集：注意在 \(\mathcal{C}(\mathrm{I})\) 中每个球都含有一个在 \([a,b[\) 上具有有界右导数的函数；另一方面，对每个 \(\varepsilon >0\) 和每个整数 \(m > 0\)，I 上存在一个连续函数，它在 \([a,b[\) 的每一点都有有限右导数，使得对一切 \(t\in [a,b[\) 有 \(|y(t)|\leqslant \varepsilon\) 与 \(|y_r'(t)|\geqslant m\)。）
\item
  设 E 是 R 上的拓扑向量空间，f 是定义在开区间 \(I \subset R\) 上的连续向量函数，并且在 I 的每一点都有右导数与左导数。
\end{enumerate}

\begin{enumerate}
\def\labelenumi{\alph{enumi})}
\item
  设 U 是 E 中的非空开集，A 是 I 中由满足 \(\mathbf{f}_{d}^{\prime}(x) \in \mathrm{U}\) 的点 x 组成的子集。给定数 \(\alpha > 0\)，设 B 是 I 中由这样的点 x 组成的子集：至少存在一个 \(y \in I\)，满足条件 \(x - \alpha \leqslant y < x\) 与 \((\mathbf{f}(x) - \mathbf{f}(y)) / (x - y) \in \mathrm{U}\)；证明集 B 是开集，并且 \(A \cap CB\) 是可数的（注意后一集由与 CB 邻接的区间的左端点组成）。由此推出，由点 \(x \in A\)（满足 \(\mathbf{f}_{g}^{\prime}(x) \notin \overline{\mathbf{U}}\)）所成的集是可数的。
\item
  假设 E 是赋范空间；这时像 \(\mathbf{f}(\mathrm{I})\) 是具有可数基的度量空间，对由 \(\mathbf{f}(\mathrm{I})\) 生成的 E 的闭向量子空间 F 亦然，该子空间包含 \(\mathbf{f}_{d}^{\prime}(\mathrm{I})\) 与 \(\mathbf{f}_{g}^{\prime}(\mathrm{I})\)。由 a) 推出，由点 \(x \in I\)（满足 \(\mathbf{f}_{d}^{\prime}(x) \neq \mathbf{f}_{g}^{\prime}(x)\)）所成的集是可数的。（若 \((\mathrm{U}_{m})\) 是 F 的拓扑的一个可数基，则注意对 F 的两个不同的点 a, b，存在两个不相交的集 \(U_{p}, U_{q}\)，使得 \(a \in U_{p}\) 与 \(b \in U_{q}\)。）
\item
  取 E 为乘积 \(R^{I}\)（即从 I 到 R 的映射所成的空间，赋予简单收敛拓扑），并对每个 \(x \in I\)，用 \(\mathbf{g}(x)\) 表示从 I 到 R 的映射 \(t \mapsto |x - t|\)。证明 g 连续，并且对每个 \(x \in I\) 有 \(\mathbf{g}_{r}'(x) \neq \mathbf{g}_{l}'(x)\)。
\end{enumerate}

\begin{enumerate}
\def\labelenumi{\arabic{enumi})}
\setcounter{enumi}{4}
\tightlist
\item
  设 \(\mathbf{f}\) 是定义在开区间 \(\mathrm{I} \subset \mathbf{R}\) 上、取值于一个赋范空间 \(\mathrm{E}\)（在 \(\mathbf{R}\) 上）中的连续向量函数，并且在 \(\mathrm{I}\) 的每一点处都有右导数。
\end{enumerate}

\begin{enumerate}
\def\labelenumi{\alph{enumi})}
\item
  证明满足如下条件的点 \(x \in \mathrm{I}\) 所成的集——\(\mathbf{f}_d'\) 在 \(x\) 的一个邻域上有界——是 \(\mathrm{I}\) 中的稠密开集（利用 Gen.~Top., IX, p.~194 的定理 2）。
\item
  证明 I 中使 \(\mathbf{f}_d^{\prime}\) 连续的点所成的集是 I 的一个瘦子集的余集（参见~Gen.~Top., IX, p.~255, 习题 21）。
\end{enumerate}

\begin{enumerate}
\def\labelenumi{\arabic{enumi})}
\setcounter{enumi}{5}
\tightlist
\item
  设 \((r_n)\) 是由 {[}0, 1{]} 中的有理数按某个次序排列而成的序列。证明函数 \(f(x) = \sum_{n=0}^{\infty} 2^{-n}(x - r_n)^{1/3}\) 连续且可导
\end{enumerate}

在 \(\mathbf{R}\) 的每个点处，并且在每个点 \(r_n\) 处有无穷导数。（为看出 \(f\) 在点 \(x\)（异于诸 \(r_n\)）处可导，分两种情形讨论，视通项为 \(2^{-n}(x - r_n)^{-2/3}\) 的级数的和是 \(+\infty\) 还是收敛而定；在第二种情形，注意对一切 \(x \neq 0\) 与一切 \(y \neq x\) 有

\[
0 \leqslant (y ^ {1 / 3} - x ^ {1 / 3}) / (y - x) \leqslant 4 / 3 x ^ {2 / 3}).
\]

\begin{enumerate}
\def\labelenumi{\arabic{enumi})}
\setcounter{enumi}{6}
\item
  设 \(f\) 是定义在区间 \(\mathrm{I} \subset \mathbf{R}\) 上的实函数，具有右导数 \(f_d'(x_0) = 0\)（在 \(\mathrm{I}\) 的一点处）；又设 \(\mathbf{g}\) 是定义在 \(y_0 = f(x_0)\) 的一个邻域上的向量函数，在该点处具有右导数与左导数（未必相等）。证明 \(\mathbf{g} \circ f\) 在点 \(x_0\) 处有等于 0 的右导数。
\item
  设 \(f\) 是从 \(\mathbf{R}\) 到其自身的映射，使得集 \(C\)（由 \(\mathbf{R}\) 中使 \(f\) 连续的点组成）在 \(\mathbf{R}\) 中稠密，并且 \(A\)（\(C\) 的余集）也稠密。证明集 \(D\)（由 \(C\) 中使 \(f\) 右可导的点组成）是瘦集。（对每个整数 \(n\)，设 \(E_n\) 是由这样的点 \(a \in \mathbf{R}\) 组成的集：存在两个点 \(x, y\)，使得 \(0 < x - a < 1/n\)、\(0 < y - a < 1/n\) 并且
\end{enumerate}

\[
\frac {f (x) - f (a)}{x - a} - \frac {f (y) - f (a)}{y - a} > 1.
\]

证明 \(E_{n}\) 的内部在 R 中稠密。为此，注意对 R 中每个非空开区间 I 都存在一点 \(b \in I \cap A\)；证明当 a \textless{} b 且 b - a 充分小时有 \(a \in E_{n}\)。

\begin{enumerate}
\def\labelenumi{\arabic{enumi})}
\setcounter{enumi}{8}
\tightlist
\item
  设 \(\mathcal{B}(\mathbf{N})\) 是实数的有界序列 \(\mathbf{x} = (x_n)_{n\in \mathbb{N}}\) 所成的空间，赋予范数 \(\| \mathbf{x}\| = \sup |x_n|\)；给出连续映射 \(t\mapsto \mathbf{f}(t) = (f_n(t))_{n\in \mathbb{N}}\)（从 \(\mathbf{R}\) 到 \(\mathcal{B}(\mathbf{N})\) 中）的一个例子，使得每个函数 \(f_{n}\) 在 \(t = 0\) 处可导，但 \(\mathbf{f}\) 在该点不可导。
\end{enumerate}

\setcounter{exercisegroup}{1}
\refstepcounter{exercisegroup}
